\section{Interupsi DOS dan Penangan Buatan Sendiri}
\label{sec:dos}

\subsection{Dua Tingkat Layanan}

BIOS menyediakan akses ke perangkat keras, tetapi belum mengenal
berkas, direktori, atau konsep ``program yang sedang berjalan''.
Semua itu adalah tugas \term{DOS} (\eng{Disk Operating System}) yang
dibangun di atas BIOS melalui interupsi \hex{21}: BIOS menggerakkan
perangkat, DOS mengatur pemakaian sistem secara keseluruhan.

\subsection{Fungsi \code{INT 21h} yang Penting}

Konvensi pemanggilan \code{INT 21h} sama dengan konvensi BIOS:
\reg{AH} berisi nomor fungsi dan register lain berisi argumen.
Fungsi yang dipakai sepanjang bab ini dirangkum pada
Tabel~\ref{tab:int21h}.

\begin{table}[H]
    \centering
    \caption{Fungsi \code{INT 21h} yang dipakai pada bab ini.}
    \label{tab:int21h}
    \small
    \begin{tabular}{@{}cp{4.1cm}p{2.5cm}p{4.7cm}@{}}
        \toprule
        \textbf{Fungsi} & \textbf{Register masukan} &
        \textbf{Register keluaran} & \textbf{Keterangan} \\
        \midrule
        \hex{01} & \reg{AH} = \hex{01} & \reg{AL} = karakter &
        Membaca satu karakter dan menampilkannya (\eng{echo});
        menunggu sampai ada tombol \\
        \addlinespace
        \hex{02} & \reg{AH} = \hex{02}; \reg{DL} = karakter & --- &
        Mencetak satu karakter pada layar \\
        \addlinespace
        \hex{09} & \reg{AH} = \hex{09}; \reg{DS}:\reg{DX} = alamat
        string & --- &
        Mencetak string yang diakhiri tanda \code{'\$'}; tanda itu
        sendiri tidak dicetak \\
        \addlinespace
        \hex{0A} & \reg{AH} = \hex{0A}; \reg{DS}:\reg{DX} = alamat
        penyangga & Penyangga terisi &
        Membaca string yang diketik pengguna sampai tombol
        \file{Enter}; karakter ditampilkan (\eng{echo}) \\
        \addlinespace
        \hex{25} & \reg{AH} = \hex{25}; \reg{AL} = nomor vektor;
        \reg{DS}:\reg{DX} = alamat penangan & --- &
        Mengganti isi entri vektor dengan alamat penangan baru \\
        \addlinespace
        \hex{35} & \reg{AH} = \hex{35}; \reg{AL} = nomor vektor &
        \reg{ES}:\reg{BX} = alamat penangan &
        Membaca isi entri vektor yang sedang berlaku \\
        \addlinespace
        \hex{4C} & \reg{AH} = \hex{4C}; \reg{AL} = kode keluar & --- &
        Mengakhiri program, menutup berkas, dan mengembalikan kendali
        ke DOS \\
        \bottomrule
    \end{tabular}
\end{table}

\begin{infobox}
Pasangan \reg{DS}:\reg{DX} adalah bentuk baku alamat jauh pada 8086
(Bab~\ref{chap:segmentasi}): \reg{DS} menyimpan nomor segmen dan
\reg{DX} menyimpan offset di dalamnya. Karena itu, program \file{.EXE}
yang memakai fungsi \hex{09} wajib menginisialisasi \reg{DS}.
\end{infobox}

\subsection{Penyangga Fungsi \hex{0A}}

Fungsi \hex{0A} berbeda dari fungsi lain karena pengguna juga
menyediakan penyangga (\eng{buffer}) yang tata letaknya ditentukan DOS
(Kode~\ref{lst:buffer-0a}).

\begin{lstlisting}[
    language=MASM,
    caption={Tata letak penyangga untuk fungsi \hex{0A}.},
    label={lst:buffer-0a}
]
.DATA
    nama    DB 20           ; byte 0: kapasitas maksimum,
                            ;   termasuk tombol Enter
            DB 0            ; byte 1: jumlah karakter yang benar-benar
                            ;   diketik, diisi oleh DOS
            DB 20 DUP(0)    ; byte 2 dan seterusnya: isi yang diketik
\end{lstlisting}

Byte pertama harus diisi \emph{sebelum} pemanggilan dan menyatakan
kapasitas maksimum, sedangkan byte kedua diisi DOS dengan jumlah
karakter yang benar-benar diketik. DOS juga menambahkan karakter
\code{0Dh}, tetapi \emph{tidak} menambahkan tanda \code{'\$'}.

Poin terakhir sering menjebak pemula: karena fungsi \hex{09} berhenti
mencetak hanya ketika menemukan tanda \code{'\$'}, program harus
menuliskan tanda itu pada posisi \code{penyangga + 2 + jumlah
karakter}, yaitu posisi \code{0Dh} yang baru saja ditulis DOS.

\begin{peringatanbox}
Bila tanda \code{'\$'} tidak dituliskan, fungsi \hex{09} akan terus
mencetak isi memori sesudah penyangga sampai menemukan byte \hex{24}
(kode ASCII \code{'\$'}). Layar pun dipenuhi karakter sampah.
Gejala ini hampir selalu berasal dari kesalahan penulisan tanda akhir
string, bukan dari kerusakan layanan DOS.
\end{peringatanbox}

\subsection{\code{INT 20h} dan \code{INT 21h} Fungsi \hex{4C}}

Dua cara mengakhiri program telah diperkenalkan: \lstasm|INT 20h| pada
Bab~\ref{chap:pengenalan} dan fungsi \hex{4C}. Perbandingannya tampak
pada Tabel~\ref{tab:terminasi}.

\begin{table}[H]
    \centering
    \caption{Perbandingan cara mengakhiri program pada DOS.}
    \label{tab:terminasi}
    \small
    \begin{tabular}{@{}p{3.4cm}p{5.2cm}p{5.2cm}@{}}
        \toprule
        \textbf{Aspek} & \textbf{\code{INT 20h}} &
        \textbf{\code{INT 21h} fungsi \hex{4C}} \\
        \midrule
        Bentuk program & Cocok untuk \file{.COM} &
        Dapat dipakai \file{.COM} maupun \file{.EXE} \\
        \addlinespace
        Prasyarat & \reg{CS} harus menunjuk ke segmen awal program &
        Tidak ada prasyarat khusus \\
        \addlinespace
        Kode keluar & Tidak dapat disertakan &
        Diberikan pada \reg{AL}, dapat dibaca program pemanggil \\
        \addlinespace
        Pembersihan oleh DOS & Terbatas &
        Menutup berkas terbuka dan memulihkan vektor
        \file{Ctrl}+\file{C} \\
        \addlinespace
        Pemakaian pada praktikum & Jarang, hanya pada contoh
        \file{.COM} & Bentuk baku program \file{.EXE} \\
        \bottomrule
    \end{tabular}
\end{table}

Kesimpulannya, fungsi \hex{4C} lebih umum dan lebih aman sehingga
menjadi pilihan baku; \lstasm|INT 20h| tetap perlu dikenal karena
muncul pada contoh program lama.

\subsection{Memasang Penangan Interupsi Sendiri}

Karena alamat penangan disimpan pada tabel vektor, mengganti isi
sebuah entri berarti mengganti seluruh perilaku interupsi tersebut.
Prosedurnya ada lima langkah: simpan penangan lama dengan fungsi
\hex{35} (hasilnya pada \reg{ES}:\reg{BX}), siapkan alamat penangan
baru pada \reg{DS}:\reg{DX}, pasang dengan fungsi \hex{25}, uji dengan
\lstasm|INT n|, lalu pulihkan penangan lama. Contoh pada
Kode~\ref{lst:pasang-vektor} memakai vektor \hex{60} yang tidak dipakai
perangkat keras mana pun sehingga aman untuk percobaan.

\begin{lstlisting}[
    language=MASM,
    caption={Menyimpan vektor lama dan memasang penangan baru.},
    label={lst:pasang-vektor}
]
    MOV AH, 35h                 ; fungsi 35h: baca vektor
    MOV AL, 60h                 ; nomor vektor 60h
    INT 21h                     ; keluaran: ES:BX = penangan lama
    PUSH ES
    PUSH BX                     ; simpan penangan lama di stack

    PUSH DS                     ; simpan DS program
    MOV DX, OFFSET handler60    ; DX = offset penangan baru
    MOV AX, SEG handler60       ; AX = segmen penangan baru
    MOV DS, AX                  ; DS:DX = alamat penangan baru
    MOV AH, 25h                 ; fungsi 25h: pasang vektor
    MOV AL, 60h
    INT 21h
    POP DS                      ; pulihkan DS program

    INT 60h                     ; uji: jalankan penangan baru
\end{lstlisting}

Penangan itu sendiri adalah prosedur jauh yang diakhiri \asm{IRET}
(Kode~\ref{lst:handler-60}). Karena dipanggil lewat interupsi dan bukan
\lstasm|CALL|, register yang dipakainya harus disimpan dan dipulihkan
sendiri.

\begin{lstlisting}[
    language=MASM,
    caption={Penangan interupsi \hex{60} yang diakhiri \asm{IRET}.},
    label={lst:handler-60}
]
handler60 PROC FAR
    PUSH AX                     ; simpan register yang akan dipakai
    PUSH DX
    MOV AH, 09h                 ; cetak pesan sapaan
    LEA DX, pesan
    INT 21h
    POP DX                      ; pulihkan register
    POP AX
    IRET                        ; pop IP, CS, lalu FLAGS
handler60 ENDP
\end{lstlisting}

Sesudah pengujian, vektor lama dipulihkan dengan mengirimkan kembali
kedua nilai yang disimpan di awal program ke fungsi \hex{25}
(Kode~\ref{lst:pulihkan-vektor}).

\begin{lstlisting}[
    language=MASM,
    caption={Memulihkan vektor interupsi yang semula.},
    label={lst:pulihkan-vektor}
]
    POP DX                      ; DX = offset penangan lama
    POP DS                      ; DS = segmen penangan lama
    MOV AH, 25h                 ; pasang kembali penangan lama
    MOV AL, 60h
    INT 21h
\end{lstlisting}

\begin{peringatanbox}
Mengganti isi vektor interupsi berisiko. Tiga kesalahan yang paling
sering terjadi: penangan tidak diakhiri \asm{IRET} sehingga prosesor
melompat ke alamat acak; vektor lama tidak dipulihkan sehingga program
berikutnya menerima penangan yang sudah tidak ada; dan penangan tidak
menyimpan register yang dipakainya. Vektor \hex{08} dan \hex{09}
sebaiknya tidak dijadikan percobaan pertama.
\end{peringatanbox}

Penangan yang menggantikan interupsi perangkat keras umumnya tidak
berdiri sendiri: sesudah pekerjaannya selesai, ia melanjutkan kendali
ke penangan lama agar fungsi asli perangkat tetap berjalan. Teknik ini
disebut \term{interupsi berantai} (\eng{interrupt chaining}) dan
dilaksanakan dengan melompat ke alamat penangan lama, bukan dengan
\asm{IRET}.
